% !TeX program = lualatex
% Compile twice: lualatex 119.tex
% All references and prose are embedded; no BibTeX database or external inputs.
\documentclass[11pt,a4paper]{article}
\usepackage[margin=24mm,headheight=14pt]{geometry}
\usepackage{fontspec}
\setmainfont{Latin Modern Roman}
\setsansfont{Latin Modern Sans}
\setmonofont{Latin Modern Mono}
\usepackage{amsmath,amssymb,mathtools}
\usepackage{unicode-math}
\setmathfont{Latin Modern Math}
\usepackage{microtype}
\usepackage{enumitem,xurl,needspace}
\usepackage[dvipsnames]{xcolor}
\usepackage{hyperref}
\hypersetup{unicode=true,colorlinks=true,linkcolor=MidnightBlue,urlcolor=MidnightBlue,
 pdftitle={500 Mathematical Problem Statements},pdfauthor={Source-annotated compilation},bookmarksnumbered=true}
\usepackage{bookmark}
\usepackage{tocloft}
\setlength{\cftsecnumwidth}{3em}
\cftsetpnumwidth{2.5em}
\cftsetrmarg{4em}
\usepackage{titlesec}
\titleformat{\section}{\Large\bfseries\raggedright}{\thesection}{0.7em}{}
\usepackage{fancyhdr}
\pagestyle{fancy}\fancyhf{}
\fancyhead[L]{\small\sffamily Mathematical problem statements}
\fancyhead[R]{\small Source edition 22}
\fancyfoot[C]{\thepage}
\setlength{\parindent}{0pt}
\setlength{\parskip}{5pt plus 1pt}
\setlength{\emergencystretch}{3em}
\setcounter{tocdepth}{1}
\setcounter{secnumdepth}{1}
\setlist[itemize]{leftmargin=3em,itemsep=5pt,topsep=4pt,parsep=0pt}
\allowdisplaybreaks
\newcommand{\N}{\mathbb{N}}
\newcommand{\Z}{\mathbb{Z}}
\newcommand{\Q}{\mathbb{Q}}
\newcommand{\R}{\mathbb{R}}
\newcommand{\C}{\mathbb{C}}
\newcommand{\F}{\mathbb{F}}
\newcommand{\A}{\mathbb{A}}
\newcommand{\PP}{\mathbb P}
\newcommand{\E}{\mathbb{E}}
\newcommand{\eps}{\varepsilon}
\newcommand{\dd}{\,\mathrm d}
\newcommand{\sref}[2]{\hyperlink{source:#1:#2}{[S#2]}}
\newcommand{\eref}[2]{\hyperlink{extra:#1:#2}{[E#2]}}
% Turn the next switch off for a shorter reading copy. The quotations preserve
% every exactTarget field, including qualifications omitted from a short summary.
\newif\ifshowcatalogue
\showcataloguetrue
\newcommand{\cataloguescope}[1]{\ifshowcatalogue
 \paragraph{Exact scope in the supplied catalogue.}
 \begin{quote}\small #1\end{quote}\fi}
\hypersetup{pdftitle={Problem 119 - Serre's congruence subgroup conjecture},pdfauthor={Open Problems Math research notebook}}
\fancyhead[L]{\small\sffamily Open Problems Math \textbar\ Research notebook}
\fancyhead[R]{\small\sffamily 119 \textbar\ 27 September 2026}
\numberwithin{equation}{section}
\begin{document}
\setcounter{section}{118}
% ================================================================
% problemId: problem.top500-omission-w041044-serre-s-conjecture-on-the-congruence-subgroup-problem
% source releaseRank: 119; source releaseStatus: open_with_solved_subcases
% exactTarget SHA-256: 59d0326282a87ede872c1a22907f60396745dd3f56efc00dabd9f638dbe9473a
\clearpage
\section{\texorpdfstring{Serre's congruence subgroup conjecture}{Serre's congruence subgroup conjecture}}\label{119}
\subsection{Definitions and mathematical statement}
For the global field $K$ and finite place set $S$, let $\mathcal O_S$ be the ring of elements integral outside $S$. The congruence completion of $G(\mathcal O_S)$ uses kernels of reduction modulo ideals, whereas its profinite completion uses all finite-index normal subgroups. Define
\[
 C^S(G)=\ker\bigl(\widehat{G(\mathcal O_S)}\to\overline{G(\mathcal O_S)}^{\rm cong}\bigr),
 \qquad r_S(G)=\sum_{v\in S}\operatorname{rank}_{K_v}G.
\]
For absolutely almost simple simply connected $G$, assert finiteness of this kernel when $r_S\ge2$ and every nonarchimedean place in $S$ has positive local rank; assert infinitude when $r_S=1$. Rank zero is excluded. Finiteness, not triviality, is the higher-rank conclusion.
\par\smallskip\textit{Formulation and concept references:} \sref{119}{1}.
\cataloguescope{Let K be a global field (a number field or the function field of a curve over a finite field), G an absolutely almost simple simply connected algebraic K-group, and S a finite nonempty set of places of K containing every archimedean place when K is a number field. Put r\_S(G) = sum over v in S of rank\_\{K\_v\}(G), and let C\textasciicircum{}S(G) be the kernel of the natural surjection from the profinite completion of G(O\_S) (S-arithmetic topology) onto its S-congruence completion. Prove Serre's congruence subgroup conjecture: (a) if r\_S(G) >= 2 and rank\_\{K\_v\}(G) > 0 for every nonarchimedean v in S, then C\textasciicircum{}S(G) is finite; (b) if r\_S(G) = 1, then C\textasciicircum{}S(G) is infinite. The case r\_S(G) = 0 (finite G(O\_S), trivial kernel) is excluded. Finiteness of C\textasciicircum{}S(G) is the assertion; triviality is a stronger statement not required. Which specific groups remain unsettled is status information, not part of the target.}
\subsection{Short English statement}
Does the congruence kernel have the finite-versus-infinite behavior predicted by the total local rank of an arithmetic group?
% BEGIN RESEARCH ADDITION 119
\subsection{Research attempt}

\paragraph{Current frontier and scope.}
The supplied survey is used for the arithmetic formulation, not as a
complete 2026 inventory of unsettled groups.  Its Sections 2.2--2.3 identify
the kernel for the arithmetic lattice with the kernel for the corresponding
completion of $G(K)$; its Theorems 1 and 2 explain how metaplectic finiteness
and arithmetic centrality enter the higher-rank problem.  It also records
infinite kernels for familiar rank-one examples, including
$\operatorname{SL}_2(\Z)$.  These are established ingredients, not a proof
for every group in the catalogue. \sref{119}{1}
The search did not provide a verified 2024--2026 result settling all the
remaining cases.  The argument below retains both global-field
characteristics and does not assume that $\Gamma=G(\mathcal O_S)$ is
finitely generated.

\paragraph{Attack routes.}
One route is to prove centrality in the arithmetic completion of $G(K)$
and apply the source's arithmetic theorems.  Its missing step is global
centrality, not the already available finiteness of the metaplectic
kernel.  A second route tests every finite quotient directly: higher rank
should bound the part invisible to congruences, whereas rank one should
produce arbitrarily large invisible parts.  A third route would build
rank-one quotients with composition factors incompatible with congruence
quotients; the construction must work for every rank-one group, not just
a virtually free example.  We develop the finite-quotient route because
it makes the required uniformity and the centrality issue explicit.

\paragraph{Attempt: exact finite defects.}
Let $q:\widehat\Gamma\twoheadrightarrow Q$ be the congruence-completion
map and $C=\ker q$.  For every finite-index normal subgroup
$N\triangleleft\Gamma$, let $U_N$ be its closure in $\widehat\Gamma$ and
set
\[
 N^{\mathrm c}=\Gamma\cap q^{-1}(q(U_N)),\qquad
 D(N)=N^{\mathrm c}/N.
\]
The symbol $\Gamma$ in this expression denotes its canonical image; the
same definition by preimage works even without first using injectivity.
The group $q(U_N)$ is compact and of finite index in $Q$, hence closed
and open.  Therefore $N^{\mathrm c}$ is a congruence-open normal subgroup
containing $N$.  It is the smallest congruence-closed subgroup containing
$N$: any such subgroup has a closed preimage containing $U_N$, hence
contains $q^{-1}(q(U_N))\cap\Gamma$.

\textbf{Finite-defect lemma.} There is a natural isomorphism
\begin{equation}\label{119:defect}
 D(N)\simeq C/(C\cap U_N).
\end{equation}
Indeed $q^{-1}(q(U_N))=CU_N$.  In the finite quotient
$\widehat\Gamma/U_N=\Gamma/N$, the image of $CU_N$ is precisely
$N^{\mathrm c}/N$.  Restricting the quotient map to $C$ proves the
claim, including surjectivity.  In particular, these quotients are not
heuristic measures of the kernel: they are its actual finite images.
If $N'\subseteq N$, the natural map $D(N')\to D(N)$ is onto.

It follows that
\begin{equation}\label{119:equivalence}
 C\text{ is finite}\quad\Longleftrightarrow\quad
 \sup_{N\triangleleft_f\Gamma}|D(N)|<\infty,
 \qquad
 |C|=\sup_N|D(N)|\quad\text{when finite}.
\end{equation}
Here $\triangleleft_f$ means normal and of finite index.  To prove the
nontrivial implication, if $C$ had more than $B$ elements, choose $B+1$
distinct ones.  The Hausdorff profinite topology separates each pair by
an open normal subgroup.  Intersect finitely many of these subgroups to
obtain some $U_N$ on which all $B+1$ images remain distinct, contradicting
the bound $B$.  This also proves the displayed equality in the finite
case by choosing a subgroup separating all elements of $C$.
No exchange of a limit and a cardinality, and no countable cofinal
sequence, is needed.

Consequently the exact two missing arithmetic assertions are
\begin{align}\label{119:targets}
 r_S\geq2\text{ with the stated local-rank condition}
 &\Longrightarrow \exists B_{K,G,S}\ \forall N:\ |D(N)|\leq B_{K,G,S},\\
 r_S=1&\Longrightarrow \forall B\ \exists N:\ |D(N)|>B.\notag
\end{align}
The constant may depend on the fixed arithmetic data.  No uniform bound
across different fields or groups has been appended.  Establishing these
assertions would prove exactly the two conclusions of the entry.

\paragraph{Testing a centrality shortcut.}
Within the lattice completion alone one has the useful equivalence
\begin{equation}\label{119:center}
 C\subseteq Z(\widehat\Gamma)
 \quad\Longleftrightarrow\quad
 [\Gamma,N^{\mathrm c}]\subseteq N\quad\text{for every }N.
\end{equation}
For the reverse implication, the image of each commutator $[c,g]$ is
trivial in every finite quotient by \eqref{119:defect}.  Finite quotients
separate points, so $[c,g]=1$ for $g\in\Gamma$.  Density and continuity
extend this to all $g\in\widehat\Gamma$.  The forward implication follows
by taking images.  Thus a failure of lattice centrality has a finite
witness.

It would be incorrect to infer \eqref{119:equivalence}'s bound merely
from \eqref{119:center}.  Consider the illustrative, non-arithmetic
pair of topologies on $\Gamma=\Z$: all finite-index subgroups versus only
$p^j\Z$ for a fixed prime $p$.  Then
\[
 \widehat{\Z}\longrightarrow\Z_p,
 \qquad C=\prod_{\ell\ne p}\Z_\ell
\]
has an infinite central kernel.  If $N=m\Z$ and
$m=p^a u$ with $(u,p)=1$, then $N^{\mathrm c}=p^a\Z$ and
$|D(N)|=u$.  Centrality holds in every finite quotient while the defects
are unbounded.  The example is not a counterexample to Serre: $\Z$ is
not an absolutely almost simple simply connected arithmetic group.
It refutes a purported general profinite lemma that would erase the
essential arithmetic step.

There is a second distinction.  The source's centrality theorem concerns
the extension of the completion of $G(K)$, not just the compact open
subgroup $\widehat\Gamma$.  Equation \eqref{119:center} does not by itself
prove that elements represented by all of $G(K)$ centralize $C$.
Applying Theorem 2 of \sref{119}{1} requires its precise ambient group
and hypotheses.  We do not silently replace them with the weaker
finite test just proved.

\paragraph{Stress test and further obstruction.}
The most attractive elementary higher-rank strategy is to write group
elements as bounded products of root elements and commute them past an
invisible finite-quotient element.  Such a proof would need a uniform
factorization theorem in the group actually at hand and compatibility
with its integral subgroups.  The catalogue includes $K$-anisotropic
groups with positive ranks at places in $S$; a supply of $K$-defined
unipotent root groups cannot be assumed there.  Thus a computation for
split elementary matrix groups would not cover the original target.
The survey distinguishes these arithmetic issues in its centrality
and generation discussion. \sref{119}{1}

Conversely, finding a noncongruence subgroup only proves $D(N)\ne1$ for
one $N$.  It does not prove an infinite kernel.  Rank one requires the
unbounded family in \eqref{119:targets}; higher rank allows a fixed
nontrivial finite kernel.  Checking finitely many quotients can therefore
certify neither conclusion universally.  The companion script checks
the $\Z$ example exactly for finite ranges of $m$, while the formula
above proves it for all $m$.

\paragraph{Outcome.}
\textbf{Outcome: Conditional reduction.}

The finite-defect description and the two equivalences
\eqref{119:equivalence}--\eqref{119:center} are proved.  They isolate the
uniform finite-quotient assertions needed for both rank regimes and expose
two invalid centrality shortcuts.  The arithmetic estimates in
\eqref{119:targets} have not been established for every permitted group.
This is a research reduction, not an extension of the known range of
Serre's conjecture, and no priority claim is made for the profinite lemmas.
% END RESEARCH ADDITION 119
\subsection{Sources}
\begin{itemize}\raggedright
\item[\textnormal{[S1]}]\hypertarget{source:119:1}{} G. Prasad and A. S. Rapinchuk, Developments on the congruence subgroup problem after the work of Bass, Milnor and Serre.
\url{https://www.mathi.uni-heidelberg.de/~wienhard/retreat10/references/prasad_rapinchuk_CSP.pdf}.
{\small\textit{Catalogue formulation source.}}
\end{itemize}


\end{document}
